\documentclass[10pt,twocolumn]{article}
\usepackage[T1]{fontenc}
\usepackage[utf8]{inputenc}
\usepackage{lmodern}
\usepackage[margin=1.8cm,top=2.4cm,bottom=2cm,columnsep=0.6cm]{geometry}
\usepackage{fancyhdr}
\usepackage{titlesec}
\usepackage{booktabs}
\usepackage{amsmath,amssymb,amsfonts}
\usepackage{microtype}
\usepackage{xcolor}
\usepackage{mdframed}
\usepackage{parskip}
\usepackage{enumitem}
\usepackage{hyperref}

\definecolor{rulered}{RGB}{180,30,30}
\definecolor{boxgray}{RGB}{245,245,245}
\definecolor{keywordgray}{RGB}{100,100,100}

\pagestyle{fancy}
\fancyhf{}
\fancyhead[L]{\textsc{\small Claude Mag}}
\fancyhead[C]{\small\textit{Machine Learning Research Digest}}
\fancyhead[R]{\small 2026-07-08}
\fancyfoot[C]{\small\thepage}
\renewcommand{\headrulewidth}{0.8pt}

\titleformat{\section}{\normalsize\bfseries\color{rulered}}{}{0pt}{}%
  [\vspace{-4pt}\color{rulered}\rule{\columnwidth}{0.4pt}]
\titlespacing{\section}{0pt}{8pt}{4pt}

\hypersetup{colorlinks=true,urlcolor=rulered,linkcolor=black}

\begin{document}
\setlength{\parindent}{0pt}
\setlength{\parskip}{4pt}

{\small\color{keywordgray}\textbf{Keywords:}
  diffusion language models \textbullet{}
  reinforcement learning \textbullet{}
  arbitrary-order generation \textbullet{}
  reasoning}\par\vspace{4pt}

{\LARGE\bfseries\raggedright
  The Flexibility Trap}\par\vspace{4pt}

{\large\itshape\raggedright
  Arbitrary-Order Generation Narrows, Not Expands, Reasoning
  in Diffusion Language Models}\par\vspace{6pt}

{\small\textbf{By} Zanlin Ni, Shenzhi Wang, Yang Yue, Tianyu Yu, Weilin Zhao,
  Yeguo Hua, Tianyi Chen, Jun Song, Cheng Yu, Bo Zheng \& Gao Huang
  (Tsinghua University et al.)}\par\vspace{2pt}

{\small\color{keywordgray}arXiv:2601.15165 \textbullet{}
  ICML 2026 Outstanding Paper Award}
\par\vspace{2pt}\noindent{\color{rulered}\rule{\columnwidth}{0.6pt}}\vspace{6pt}

Diffusion large language models promise to supersede autoregressive generation by
producing tokens in any order, not just left-to-right. The promise has a flaw:
for reasoning tasks like mathematics and coding, that same flexibility teaches
models to \emph{avoid} uncertain tokens rather than resolve them, collapsing the
solution space before it has been explored. A minimalist fix---JustGRPO,
which constrains rollouts to left-to-right during RL training only---achieves
\textbf{89.1\%} on GSM8K, outperforming every prior diffusion RL method by more
than six percentage points.

\begin{mdframed}[backgroundcolor=boxgray,linecolor=rulered,linewidth=1pt,
    innerleftmargin=6pt,innerrightmargin=6pt,
    innertopmargin=6pt,innerbottommargin=6pt]
\textbf{\small AT A GLANCE}\par\vspace{4pt}
\begin{description}[leftmargin=0pt,labelsep=4pt,itemsep=2pt,parsep=0pt,
    font=\small\bfseries]
  \item[Problem:] Diffusion LLMs exploit arbitrary-order generation to skip
    high-uncertainty tokens during RL exploration, collapsing solution coverage.
  \item[Why it matters:] dLLMs are positioned as the next generation of
    language model architectures; if their defining feature sabotages reasoning,
    that claim requires fundamental reassessment.
  \item[Method:] Coverage analysis of rollout distributions; JustGRPO constrains
    exploration to left-to-right ordering while leaving inference unchanged.
  \item[Result:] 89.1\% on GSM8K and 45.1\% on MATH-500, surpassing best
    prior diffusion RL methods (82.8\% and 39.6\%) by large margins.
  \item[Evidence:] Comparisons against ESPO and GDPO on GSM8K, MATH-500, and
    AIME; ablations confirm the constraint matters only at training time.
\end{description}
\end{mdframed}

\section{The Big Picture}

Diffusion language models (dLLMs) generate text by iteratively denoising a
masked token sequence, selecting which token to reveal at each step in any
order. This departs from the strictly left-to-right generation of autoregressive
transformers and, in theory, should allow dLLMs to explore a solution space that
strictly subsumes the autoregressive one: any left-to-right trajectory is a
special case of arbitrary-order generation.

Recent work applied reinforcement learning to dLLMs for reasoning tasks,
producing methods like ESPO and GDPO. These inherit the arbitrary-order
assumption uncritically. This paper asks a prior question: \emph{is that
assumption correct?}

The answer is no. Arbitrary-order generation, when combined with RL exploration,
traps models in a degenerate strategy that undermines the very diversity RL is
meant to exploit.

\section{Why Existing Approaches Fall Short}

ESPO and GDPO apply group relative policy optimization directly to dLLM rollouts,
allowing arbitrary generation order throughout. The coverage analysis reveals
the consequence: models learn to generate high-confidence (low-entropy) tokens
first, committing early to a particular solution path, and defer high-uncertainty
tokens to the end where their influence on the output is minimal.

This is structurally equivalent to a tree search algorithm that always expands
its most confident node and never backtracks. The diversity in the rollout
batch that GRPO needs for advantage estimation collapses: many rollouts
converge to the same prefix, and the effective group size shrinks.

Prior interpretations attributed performance gaps to model size or training
duration. The coverage diagnostic shows the bottleneck is the order constraint:
the problem is present from the first RL iteration.

\section{Core Mechanism}

JustGRPO makes a single modification to the dLLM RL pipeline: during the
rollout generation phase, tokens are produced strictly left-to-right. The
model's bidirectional attention and parallel decoding at inference are
completely unchanged.

This forces the model to resolve uncertain tokens in context order during
exploration, maintaining a diverse rollout distribution. The GRPO advantage
estimator then has access to genuine diversity---rollouts that took genuinely
different paths---producing meaningful gradient signal.

At inference, the model reverts to its native arbitrary-order parallel decoding,
so the speed benefits of dLLMs are fully preserved.

\section{Technical Formulation}

Let $G$ rollouts $\{o_1, \ldots, o_G\}$ be sampled from the current policy
$\pi_\theta$ for prompt $q$. Define normalized advantages
\[
\hat{A}_i = \frac{r_i - \text{mean}(r_{1:G})}{\text{std}(r_{1:G})}
\]
where $r_i \in \{0, 1\}$ is the binary correctness reward. The GRPO objective is
\[
\mathcal{J}(\theta) = \mathbb{E}\!\left[\frac{1}{G}\sum_{i=1}^{G}\frac{1}{|o_i|}
  \sum_{t=1}^{|o_i|}\min\!\left(
    \hat{\rho}_{i,t}\hat{A}_i,\,
    \text{clip}\!\left(\hat{\rho}_{i,t},1{-}\epsilon,1{+}\epsilon\right)\!\hat{A}_i
  \right)\!\right]
  - \beta D_{\mathrm{KL}}(\pi_\theta \| \pi_{\mathrm{ref}})
\]
where $\hat{\rho}_{i,t} = \pi_\theta(o_{i,t}\mid q, o_{i,<t}) / \pi_{\mathrm{ref}}(o_{i,t}\mid q, o_{i,<t})$.
JustGRPO differs only in that rollouts $o_i$ are generated left-to-right.

Coverage collapse is measured by the expected entropy of the rollout prefix
distribution at each position $k$:
\[
\mathcal{H}_k = H\!\left(\{o_{i,1:k}\}_{i=1}^{G}\right).
\]
Under arbitrary-order dLLM rollouts, $\mathcal{H}_k$ drops sharply by position 3;
under JustGRPO it remains high throughout.

\section{Learning or Inference Procedure}

\begin{enumerate}[itemsep=2pt,parsep=0pt]
  \item Sample a math/code problem $q$ from the training set.
  \item Generate $G = 8$ rollouts using \textbf{left-to-right decoding} with
    temperature 1.0.
  \item Score each rollout for correctness: $r_i \in \{0,1\}$.
  \item Compute GRPO advantages $\hat{A}_i$ from the group reward distribution.
  \item Update $\theta$ via the clipped surrogate loss with KL penalty $\beta$.
  \item At \textbf{inference}, use native parallel decoding (arbitrary order).
\end{enumerate}

\section{What the Guarantee Says}

The paper does not provide formal convergence guarantees, but proves
empirically that coverage collapse is absent under left-to-right constraint
and universal under arbitrary-order generation across all tested model families.
The key invariant is: $\mathcal{H}_k^{\text{JustGRPO}} \geq c \cdot k$ for
a positive constant $c$ throughout training, while
$\mathcal{H}_k^{\text{arbitrary}} \to 0$ within a few gradient steps.

\section{Experimental Findings}

\begin{center}
\small
\begin{tabular}{lcc}
\toprule
\textbf{Method} & \textbf{GSM8K} & \textbf{MATH-500} \\
\midrule
ESPO   & 82.3\% & 39.0\% \\
GDPO   & 82.8\% & 39.6\% \\
\textbf{JustGRPO} & \textbf{89.1\%} & \textbf{45.1\%} \\
\bottomrule
\end{tabular}
\end{center}

On AIME 2024 (competition-level math), JustGRPO reaches 23.3\%
vs.\ GDPO's 17.8\%. Closed-source dLLM baselines without RL lag further behind.
JustGRPO uses the same base model (MDLM) and identical training compute as the
comparison methods.

\section{Ablations and Interpretation}

\textbf{Order constraint at inference:} Forcing left-to-right at inference
\emph{reduces} accuracy by 1.8 points, confirming the constraint should
apply to training only.

\textbf{Group size $G$:} Performance is robust to $G \in \{4, 8, 16\}$;
the coverage collapse under arbitrary order worsens with larger $G$ because
degenerate prefixes dominate faster.

\textbf{Base model:} Results replicate on MDLM, PLAID, and a dLLM based on
LLaMA-3-8B. The effect is architectural, not model-specific.

\textbf{Task type:} Creative writing shows minimal gap (0.3 points),
confirming the effect is specific to tasks with verifiable, structured
correct answers where high-uncertainty tokens are decision-critical.

\section*{Reference}

Zanlin Ni, Shenzhi Wang, Yang Yue, Tianyu Yu, Weilin Zhao, Yeguo Hua,
Tianyi Chen, Jun Song, Cheng Yu, Bo Zheng, Gao Huang.
\textbf{The Flexibility Trap: Rethinking the Value of Arbitrary Order in
Diffusion Language Models}.
arXiv:2601.15165, January 2026.
\url{https://arxiv.org/abs/2601.15165}

\end{document}
